\answer{ex:15:1}{(a)~$2\cdot10^5$ par neurone, $6\cdot10^{12}$ pour le circuit. (b)~$\approx8$~bit/s par fil, au moins $2.5\cdot10^4$~s $\approx7$~h.}
\answer{ex:15:2}{Facteur $2$: corrélation entre voisins $0.943$, $\lambda$ de $4.18$ à $7.3\cdot10^{-4}$, rapport $\approx5.7\cdot10^3$; facteur $4$: corrélation $0.8$, rapport $\approx94$.}
\answer{ex:15:3}{(a)~$\lambda=0.988$ et $0.808$: $82$ et $4.7$ mouvements. (b)~$x_s^*=0.690$, $x_f^*=0.172$, somme $0.862$.}
\answer{ex:15:4}{(a)~$116$ mouvements. (b)~$19$. Un modèle à un seul processus, à somme nulle, n'apporte aucune économie.}
\answer{ex:15:5}{(a)~$G=50$~s$^{-1}$, contre une limite de $10.5$~s$^{-1}$ sans modèle. (b)~Non: pour $G=50$~s$^{-1}$, le retard critique est $d_c\approx103\ms$; pour $d=150\ms$, il faut $\hat k>0.445k$ (pour tous $G$ et $d$, $\hat k>k/2$).}
\answer{ex:15:6}{L'erreur passe sous $0.5\%$ en $6$--$30$~s, avec un plancher de $\approx(6$--$9)\cdot10^{-4}$ pour $7.5\cdot10^{-4}$ avec les meilleurs poids; pour $\eta=50$, les poids diffèrent encore des meilleurs de jusqu'à $0.2$ le long des directions mal conditionnées de $C$ (exercice~\ref{ex:15:2}).}
\answer{ex:15:7}{$\mathbf B=A\mathbf K$; $q_f/R\approx0.035$, $q_s/R\approx8.6\cdot10^{-4}$; avec $4q_s$, $B_f\approx0.088$, $B_s\approx0.047$: le processus lent apprend plus de deux fois plus vite, le rapide plus lentement.}
